\chapter{BNT and Full Assembly}\label{ch:assembly}

This chapter introduces \emph{Basis Normal Tensors} (BNT),
establishes their permutation rigidity, and records the
strongest end-stage Fundamental Theorem statements currently
assembled in Lean.
The proportional theorem is formalized with explicit
coefficient-convergence data, and the equal theorem is the
same-structure CF-BNT specialization rather than the full
paper-level equal-MPV theorem.
The results combine the canonical form reduction
(Chapter~\ref{ch:canonical}),
block separation (Chapter~\ref{ch:block_perm}),
spectral gap (Chapter~\ref{ch:spectral}),
and BNT permutation arguments; the Newton--Girard identities
are included as supporting infrastructure for future
weight-matching extensions.
The main references are
\cite{PerezGarcia2007Matrix} and~\cite{Cirac2021Matrix}.

\section{Basis normal tensors}

\begin{definition}[Basis Normal Tensor (BNT)]\label{def:bnt}
	\lean{MPSTensor.IsBNT}
	\leanok
	\uses{def:normal, def:mpv}
	A \emph{basis normal tensor} (BNT) decomposition for a
	total tensor $A_{\mathrm{tot}}$ consists of $g$ block tensors
	$(A_1, \ldots, A_g)$ with bond dimensions $(D_1, \ldots, D_g)$
	such that:
	\begin{enumerate}
		\item each $A_j$ is normal
		      (Definition~\ref{def:normal}),
		\item at each system size~$N$, the MPV of $A_{\mathrm{tot}}$
		      is a linear combination of the block MPVs:
		      $\ket{V^{(N)}(A_{\mathrm{tot}})}
			      = \sum_j c_j \ket{V^{(N)}(A_j)}$,
		\item for sufficiently large~$N$, the block MPV
		      vectors $\{\ket{V^{(N)}(A_j)}\}_{j=1}^g$ are
		      linearly independent.
	\end{enumerate}
	This is \cite[Definition~4.2]{Cirac2021Matrix}.
\end{definition}

\begin{definition}[Canonical form with BNT separation]\label{def:cf_bnt}
	\lean{MPSTensor.IsCanonicalFormBNT}
	\leanok
	\uses{def:is_canonical_form, def:gauge_phase_equiv}
	A \emph{canonical form with BNT separation} extends the canonical form
	(Definition~\ref{def:is_canonical_form})
	with the additional requirement that distinct blocks are
	\emph{not} gauge-phase equivalent:
	for $j \neq k$ with $D_j = D_k$,
	$A_j$ and $A_k$ are not gauge-phase equivalent.
\end{definition}

\begin{theorem}[Cross-overlap decay for distinct CF-BNT blocks]\label{thm:cross_overlap_zero}
	\lean{MPSTensor.IsCanonicalFormBNT.cross_overlap_tendsto_zero}
	\leanok
	\uses{def:cf_bnt, def:mpv_overlap}
	If $((\mu_k),(A_k))$ satisfies the CF-BNT predicate and
	$j \neq k$, then $O_{A_j A_k}(N) \to 0$.
\end{theorem}

\begin{proof}\leanok\uses{thm:overlap_decay, thm:overlap_decay_rect}
	If $D_j \neq D_k$, the rectangular spectral gap theorem
	(Theorem~\ref{thm:overlap_decay_rect}) gives
	$O_{A_j A_k}(N) \to 0$.
	If $D_j = D_k$, the blocks are not gauge-phase
	equivalent (by the CF-BNT separation condition),
	so the spectral-gap theorem
	(Theorem~\ref{thm:overlap_decay}) gives decay to~$0$.
\end{proof}

\begin{theorem}[CF-BNT yields BNT]\label{thm:cf_bnt_is_bnt}
	\lean{MPSTensor.IsCanonicalFormBNT.isBNT}
	\leanok
	\uses{def:cf_bnt, def:bnt}
	If $((\mu_k), (A_k))$ satisfies the CF-BNT predicate, then
	the block tensors form a BNT decomposition for the
	block-diagonal tensor $\bigoplus_k \mu_k A_k$.
\end{theorem}

\begin{proof}
	\leanok
	\uses{thm:mpv_decomposition, thm:cross_overlap_zero}
	\textbf{Normality}: each block is injective
	(hence $1$-block injective, hence normal).
	\textbf{Spanning}: the MPV decomposition
	(Theorem~\ref{thm:mpv_decomposition}) gives
	$V^{(N)}(A_{\mathrm{tot}})_\sigma = \sum_k \mu_k^N V^{(N)}(A_k)_\sigma$.
	\textbf{Linear independence}: the overlap matrix
	$G_{jk}(N) = O_{A_j A_k}(N)$ converges to the identity
	(diagonal entries $\to 1$ by the self-overlap hypothesis;
	off-diagonal entries $\to 0$ by
	Theorem~\ref{thm:cross_overlap_zero}).
	A Gram matrix converging to the identity is eventually
	invertible, giving linear independence.
\end{proof}

\section{BNT permutation rigidity}

\begin{theorem}[BNT permutation --- proportional MPV case]\label{thm:bnt_perm_prop}
	\lean{MPSTensor.exists_eq_numBlocks_and_equiv_gaugePhase_of_proportional_decomp}
	\leanok
	\uses{def:bnt, def:gauge_phase_equiv}
	Let $(A_j)_{j=1}^{g_A}$ and $(B_k)_{k=1}^{g_B}$ be two
	BNT decompositions such that the block-diagonal tensors
	generate proportional MPV families with convergent nonzero
	coefficients.
	Then $g_A = g_B$, and there is a permutation
	$\sigma \in S_g$ with $D_{\sigma(k)} = D_k$ and
	$A_{\sigma(k)}$ gauge-phase equivalent to $B_k$ for each~$k$.
\end{theorem}

\begin{proof}
	\leanok
	From the proportional MPV hypothesis and convergent
	coefficients, one extracts (for each pair $(j,k)$) the
	limit $\lim_{N \to \infty} O_{A_j B_k}(N)$.
	By the spectral gap (Chapter~\ref{ch:spectral}),
	cross-block overlaps between non-equivalent blocks
	decay to~$0$.
	The resulting overlap matrix has at most one nonzero entry
	per row and per column, which defines a permutation.
	Nonzero overlap forces gauge-phase equivalence.

	This is part of \cite[Theorem~4.1]{Cirac2021Matrix}
	(the Fundamental Theorem for proportional MPVs),
	which establishes the permutation and gauge-phase
	equivalence from the uniqueness of BNT decompositions
	(characterised in \cite[Proposition~4.1]{Cirac2021Matrix}).
\end{proof}

\begin{theorem}[BNT permutation --- overlap-orthonormal case]\label{thm:bnt_perm_simple}
	\lean{MPSTensor.exists_eq_numBlocks_and_equiv_gaugePhase_of_overlapOrtho}
	\leanok
	\uses{def:bnt, def:gauge_phase_equiv}
	Let $(A_j)_{j=1}^{g_A}$ and $(B_k)_{k=1}^{g_B}$ be two
	families of injective DS-gauge tensors whose self-overlaps
	converge to~$1$ and whose cross-overlaps (within each
	family) converge to~$0$.
	If they span the same MPV subspace at every system size,
	then $g_A = g_B$ and the blocks are related by a
	permutation and gauge-phase equivalence.
\end{theorem}

\begin{proof}\leanok\uses{thm:overlap_decay}
	The overlap matrix $O_{A_j B_k}(N)$ converges to a
	matrix with at most one nonzero entry per row/column
	(by the spectral gap).
	Thus the overlap orthonormality hypotheses are
	assumed directly, rather than derived from
	proportional MPVs together with coefficient convergence.
\end{proof}

\section{Coefficient convergence}

\begin{theorem}[Coefficient ratio decay]\label{thm:coeff_ratio_decay}
	\lean{MPSTensor.IsCanonicalForm.coeff_ratio_tendsto}
	\leanok
	\uses{def:is_canonical_form}
	Under the canonical form predicate, for $k \ge 1$,
	$(\mu_k / \mu_0)^N \to 0$
	as $N \to \infty$.
\end{theorem}

\begin{proof}\leanok
	Since $|\mu_k| < |\mu_0|$ (strict ordering of moduli),
	$|\mu_k / \mu_0| < 1$, so the geometric sequence
	converges to~$0$.
\end{proof}

\begin{remark}[Strict-dominance regime]\notready
	Theorem~\ref{thm:coeff_ratio_decay} applies when one
	modulus is strictly dominant:
	$|\mu_0| > |\mu_k|$ for $k \ge 1$.
	The oscillatory case
	$\sum_q \mu_{j,q}^N$ with $|\mu_{j,q}| = 1$
	is treated separately in the full paper and is not formalized here.
\end{remark}

\section{Newton--Girard identities}

\begin{theorem}[Newton--Girard trace recursion]\label{thm:newton_girard}
	\lean{Matrix.charpoly_eq_of_forall_trace_pow_eq}
	\leanok
	If $A, B \in \MN{n}$ satisfy
	$\tr(A^k) = \tr(B^k)$ for all $k \ge 1$,
	then $A$ and $B$ have the same characteristic polynomial.
\end{theorem}

\begin{proof}
	\leanok
	The Newton--Girard identities express the coefficients
	of the characteristic polynomial recursively in terms
	of the power-sum traces $\tr(A^k)$.
	Equal power sums give equal coefficients, hence equal
	characteristic polynomials.
\end{proof}

\begin{theorem}[Equal power sums imply equal multisets]\label{thm:power_sum_multiset}
	\lean{Matrix.sum_pow_eq_implies_multiset_eq}
	\leanok
	Let $\alpha, \beta : \{1, \ldots, n\} \to \C$.
	If $\sum_i \alpha_i^k = \sum_i \beta_i^k$ for all $k \ge 1$,
	then the multisets $\{\alpha_i\}$ and $\{\beta_i\}$ are equal.
\end{theorem}

\begin{proof}
	\leanok
	\uses{thm:newton_girard}
	Construct diagonal matrices $\diag(\alpha)$ and
	$\diag(\beta)$.
	The power-sum hypothesis gives
	$\tr(\diag(\alpha)^k) = \tr(\diag(\beta)^k)$.
	By Theorem~\ref{thm:newton_girard}, the characteristic
	polynomials agree: $\prod_i (X - \alpha_i) = \prod_i (X - \beta_i)$.
	Extracting roots gives multiset equality.
\end{proof}

\section{Primitivity bridge}

\begin{definition}[Primitive MPS tensor]\label{def:primitive_mps}
	\lean{MPSTensor.IsPrimitiveMPS}
	\leanok
	\uses{def:mps_tensor}
	A \emph{primitive MPS tensor} is one equipped with a nonzero
	PSD fixed point $\rho$ of the transfer map such that the
	complementary spectral radius
	$\rho_{\mathrm{spec}}(\E_A - P) < 1$
	(where $P$ is the fixed-point projection).
\end{definition}

\begin{theorem}[Primitive overlap convergence]\label{thm:prim_overlap_one}
	\lean{MPSTensor.IsPrimitiveMPS.overlap_tendsto_one}
	\leanok
	\uses{def:primitive_mps, def:mpv_overlap}
	If $A$ is a primitive MPS tensor, then
	$O_{AA}(N) \to 1$.
\end{theorem}

\begin{proof}\leanok\uses{thm:primitive_overlap}
	Immediate from the primitive overlap convergence
	(Theorem~\ref{thm:primitive_overlap}).
\end{proof}

\section{Full assembly}

\begin{theorem}[Fundamental Theorem --- proportional MPV case]\label{thm:ft_proportional}
	\lean{MPSTensor.fundamentalTheorem_proportionalMPV_CFBNT}
	\leanok
	\uses{def:cf_bnt, def:gauge_phase_equiv}
	Let $A_{\mathrm{tot}}$ and $B_{\mathrm{tot}}$ be MPS tensors
	equipped with CF-BNT families $(\mu_j^A, A_j)_{j=1}^{r_A}$ and
	$(\mu_k^B, B_k)_{k=1}^{r_B}$.
	Assume one is given coefficient arrays $a_{N,j}$, $b_{N,k}$ and
	nonzero limits $a_j^\infty$, $b_k^\infty$, $c^\infty$ such that
	\[
		V^{(N)}(A_{\mathrm{tot}})_\sigma
		= \sum_{j=1}^{r_A} a_{N,j} V^{(N)}(A_j)_\sigma,
		\qquad
		V^{(N)}(B_{\mathrm{tot}})_\sigma
		= \sum_{k=1}^{r_B} b_{N,k} V^{(N)}(B_k)_\sigma,
	\]
	with $a_{N,j} \to a_j^\infty \neq 0$ and
	$b_{N,k} \to b_k^\infty \neq 0$, and suppose also that
	\[
		V^{(N)}(A_{\mathrm{tot}})_\sigma
		= c_N V^{(N)}(B_{\mathrm{tot}})_\sigma
	\]
	for a sequence $c_N \to c^\infty \neq 0$.
	Then $r_A = r_B$, and there is a permutation $\pi$ such that
	$D_j^A = D_{\pi(j)}^B$ and $A_j$ is gauge-phase equivalent to
	$B_{\pi(j)}$ for each~$j$.
	The coefficient and convergence data are hypotheses of the Lean
	theorem; they are not yet derived automatically from arbitrary
	input tensors.
\end{theorem}

\begin{proof}
	\leanok
	\uses{thm:cf_bnt_is_bnt, thm:bnt_perm_prop}
	The CF-BNT hypotheses give BNT decompositions by
	Theorem~\ref{thm:cf_bnt_is_bnt}.
	With the supplied decomposition identities and convergence
	hypotheses on the coefficients and proportionality constants,
	Theorem~\ref{thm:bnt_perm_prop} yields the equality of block
	counts, the permutation, and the per-block gauge-phase
	equivalence.
\end{proof}

\begin{theorem}[Fundamental Theorem --- equal MPV case]\label{thm:ft_equal}
	\lean{MPSTensor.fundamentalTheorem_equalMPV_CFBNT}
	\leanok
	\uses{def:cf_bnt, def:gauge_equiv}
	Fix a common block structure: the same number of blocks~$r$,
	the same bond dimensions $(D_k)_{k=1}^r$, and the same weights
	$(\mu_k)_{k=1}^r$.
	Let $(\mu_k, A_k)$ and $(\mu_k, B_k)$ be two CF-BNT families on
	this common data.
	If the block-diagonal tensors
	$\bigoplus_k \mu_k A_k$ and $\bigoplus_k \mu_k B_k$
	generate the same MPV family, then each pair $(A_k, B_k)$ is
	gauge equivalent, and hence the assembled block-diagonal tensors
	are gauge equivalent.
	This is the same-structure equal-MPV specialization formalized in
	Lean, not the full paper-level equal-MPV theorem with
	independently given block data.
\end{theorem}

\begin{proof}
	\leanok
	\uses{thm:ft_canonical}
	Forget the extra BNT-separation field and retain only the
	underlying canonical-form data.
	The canonical-form fundamental theorem
	(Theorem~\ref{thm:ft_canonical}) then applies directly to this
	common $(r, D_k, \mu_k)$ structure, yielding per-block gauge
	equivalence and hence global gauge equivalence.
\end{proof}
